\documentclass[10pt,twocolumn,letterpaper]{article}

% Standalone one-page evidence-table workspace. It deliberately retains the
% official ACCV rebuttal geometry and typography from rebuttal copy.tex.
\usepackage[rebuttal]{cvpr}
\usepackage{accvabbrv}
\usepackage{graphicx}
\usepackage{booktabs}
\usepackage{array}
\definecolor{accvblue}{rgb}{0.12,0.49,0.85}
\usepackage[pagebackref,breaklinks,colorlinks,citecolor=accvblue]{hyperref}

\def\paperID{1107}
\def\confName{ACCV}
\def\confYear{2026}
\setcounter{table}{5} % Continue after Tables 1--5 in the submission.

\begin{document}
\title{Exploring RGB and Video V-JEPA 2.1 Features for Aerial Depth Estimation and Semantic Scene Completion}
\maketitle
\thispagestyle{empty}
\appendix

\noindent\textbf{We thank all reviewers for their constructive feedback.}
VJ/D2/D3/DA2/MVSF abbreviate V-JEPA~2.1/DINOv2/DINOv3/Depth Anything~2/MVSFormer++;
response table numbering starts at 6, continuing submission.
\begin{table}[t]
  \centering
  \caption{\textbf{Evaluation of frozen foundation-model features with increasingly capable monocular depth heads and posed multi-view stereo integration.}}
  \label{tab:depth_evaluation}
  \footnotesize
  \setlength{\tabcolsep}{2.5pt}
  \begin{tabular}{@{}lcccc@{}}
    \toprule
    Method & $\delta_1\uparrow$ & AbsRel$\downarrow$ & RMSE$\downarrow$ & SILog$\downarrow$ \\
    \midrule
    VJ DPT          & 0.747 & 0.162 & 6.730 & 0.200 \\
    VJ video Linear Probe              & 0.701 & 0.191 & 7.055 & 0.207 \\
    D3 Linear Probe                    & 0.782 & 0.156 & 6.097 & 0.186 \\
    D3 FiLM-DPT              & \textbf{0.848} & \textbf{0.121} & \textbf{4.358} & 0.161 \\
    VJ DTA-MVS (pending)     & -- & -- & -- & -- \\
    VJ DTA-FiLM-MVS          & 0.768 & 0.175 & 6.406 & 0.211 \\
    MVSF++/D2-B             & 0.748 & 0.231 & 7.657 & 0.195 \\
    MVSF++/VJ-G             & 0.761 & 0.225 & 7.517 & 0.193 \\
    MVSF++/D3-H+            & 0.828 & 0.136 & 4.611 & \textbf{0.130} \\
    \bottomrule
  \end{tabular}
\end{table}

\paragraph{Responses common to LD5V, CWVJ, and y7qY:}
 Table~\ref{tab:depth_evaluation} caption is appropriate for the Table 3 and it extends Table 3 by reporting corrected result for DPT, Spatial alignment included video methods, suggested DINO variants for faithful comparison and attribution of merits or flaws, requested MVS(simple) and supportive VJ/D features backed MVSF++(sofisticated) variants depth results.  Only important baseline is DA2-Occufly. For other baselines refer table 5 of occufly paper. We omit superiority claim of VJ FiLM-DPT. Considering the Table~3 and Table~\ref{tab:depth_evaluation} we show increasingly capable depth heads and MVS integrations. Although DAv2-OccuFly achieves the strongest reported SILog and threshold accuracies, foundation models with general pretraining objectives remain competitive and outperform it on absolute-error metrics. Any remaining gap may reflect DAv2’s depth-specific pretraining and end-to-end OccuFly adaptation, whereas frozen VJ/D backbones are adapted only through their depth heads.

\begin{table}[t]
  \caption{\textbf{Altitude interpolation diagnostics.} }
  \label{tab:altitude-interpolation-control}
  \centering
  \footnotesize
  \setlength{\tabcolsep}{2.2pt}
  \resizebox{\columnwidth}{!}{%
  \begin{tabular}{@{}lccccc@{}}
    \toprule
    Method & $\delta_1\uparrow$ & AbsRel$\downarrow$ & RMSE$\downarrow$ &
      MAE$\downarrow$ & SILog$\downarrow$ \\
    \midrule
    \multicolumn{6}{@{}l}{\textbf{A. Unseen altitude; familiar scenes}} \\
    \multicolumn{6}{@{}l}{Tr: 01--08@30/50\,m; Val: 01--04@40\,m; Te: 05--08@40\,m} \\
    DPT 
      & 0.789 & 0.151 & 6.728 & 5.775 & 0.180 \\
    FiLM-DPT
      & \textbf{0.965} & \textbf{0.102} & \textbf{4.707} &
        \textbf{3.957} & \textbf{0.120} \\
    \midrule
    \multicolumn{6}{@{}l}{\textbf{B. Unseen scenes and altitude}} \\
    \multicolumn{6}{@{}l}{Tr: 01--04@30/50\,m; Val: 05--06@40\,m; Te: 07--08@40\,m} \\
    DPT
      & \textbf{0.925} & \textbf{0.102} & \textbf{4.807} &
        \textbf{3.994} & \textbf{0.125} \\
    FiLM-DPT
      & 0.707 & 0.205 & 8.506 & 7.918 & 0.169 \\
    \bottomrule
  \end{tabular}
  }
\end{table}

Table~\ref{tab:altitude-interpolation-control} shall be Table~4. The main purpose of this non standard split ablation is to specifically study methods ability to interpolate to unseen altitude cues. But as rigltly pointed there is inherently visual cues leakage. To address that we will report a ablation of unseen scenes and altitude for comparision.
A clear architecture diagram for SSC concept, prioritizing data-flow and various components involved can replace Figures~1 and~5 in the final submission.


\begin{table}[t]
  \caption{\textbf{FiLM conditioningablation.} Test mean$\pm$SD
  over seeds 0/1/2. Numerically best mean in \textbf{bold}.}
  \label{tab:film-aggregate}
  \centering
  \footnotesize
  \setlength{\tabcolsep}{2.6pt}
  \resizebox{\columnwidth}{!}{%
  \begin{tabular}{@{}lcccc@{}}
    \toprule
    Conditioning & $\delta_1\uparrow$ & AbsRel$\downarrow$ & RMSE$\downarrow$ & SILog$\downarrow$ \\
    \midrule
    DPT
      & $.728{\pm}.034$ & $.174{\pm}.006$ & $6.709{\pm}.385$ & $.202{\pm}.010$ \\
    Altitude FiLM
      & $.798{\pm}.020$ & $.145{\pm}.005$ & $5.236{\pm}.322$ & $.180{\pm}.006$ \\
    Intrinsics FiLM
      & $.711{\pm}.015$ & $.182{\pm}.004$ & $7.115{\pm}.243$ & $.216{\pm}.007$ \\
    Alt.+intr. FiLM
      & $\mathbf{.810{\pm}.006}$ & $\mathbf{.137{\pm}.005}$ & $\mathbf{4.907{\pm}.140}$ & $\mathbf{.171{\pm}.005}$ \\
    \bottomrule
  \end{tabular}
  }
\end{table}
\begin{table}[t]
  \caption{\textbf{Extension of main-paper Table~5 with MVSFormer++
  depth into FSSC geometric layer.} }
  \label{tab:fssc-mvsformerpp-extension}
  \centering
  \footnotesize
  \setlength{\tabcolsep}{3.2pt}
  \resizebox{\columnwidth}{!}{%
  \begin{tabular}{@{}lcc@{}}
    \toprule
    Variant & SC IoU (\%)$\uparrow$ & SSC mIoU (\%)$\uparrow$ \\
    \midrule
    DINOv3 context + DINOv3 MVS depth
      & \textbf{43.37} & \textbf{4.47} \\
    DINOv3 context + DINOv3 MVS depth + VoxDet V5-V2
      & \textbf{48.06} & \textbf{4.84} \\
    V-JEPA2.1 context + V-JEPA2.1 MVS depth
      & 35.02 & 3.28 \\
    \bottomrule
  \end{tabular}
  }
\end{table}

\paragraph{LD5V:statistical analysis of
robustness, and evidence for the central claims.}
contribution of the FiLM ablation: We provide a matched FiLM ablation and its per-seed results in Table~\ref{tab:film-aggregate} and Table~\ref{tab:film-seeds}.
spatial misalignment: We thank the reviewer for pointing this out and adopted 2D feature alignment like frozen RAFT flow to align neighboring features to the target frame;
Alongside simple MVS baselines, extending with MVSFormer++, a sophisticated architecture combining cross-view fusion with attention-based feature encoding and cost-volume regularization for depth estimation, is appropriate. Ablations replacing the original D2 backbone with VJ and D3 for depth and FSSC-style geometric lifting are reported in Table~\ref{tab:fssc-mvsformerpp-extension}. Adding the VoxDet V5-V2 recipe (full-resolution VoxNT, long-tail loss, and scheduled depth gradients) to the D3 MVSFormer++ variant reaches 48.06 SC IoU and 4.84 union-present SSC mIoU; its GT-present SSC mIoU is 5.11.
Oracle depth:Replacing estimated depth with GT metric depth in the same geometric lifting
and SSC pipeline yields SC IoU/SSC mIoU of 49.40/7.25 with V-JEPA2.1 ViT-G RGB
context and 42.01/8.40 with DINOv3 ViT-H/16+ RGB context.
per-class results: Already provided in the supplemnetary material section B. Will be further extended with MVSF++FSSC variants.

\paragraph{CWVJ:missing cost figures and related work, existing
aerial datasets and multi-dataset evaluation, framing the conclusion as a
study rather than a new SSC architecture, the confounded key baseline, and
evaluation on only one benchmark.}
DINOv3: We provide an ablation table containing DINOv3 alongside V-JEPA~2.1.
This important point also motivates, if permitted, a revised title:
``Exploring Foundation-Model Spatial and Temporal Features for Aerial Depth
Estimation and Semantic Scene Completion.''
Single-run:Table \ref{film-aggregate} for FiLM contribution ablations given the
limited rebuttal period, and can extend this analysis to all published results.
Because this preliminary study examines frozen foundation features for MDE and SSC rather than onboard deployment, we defer SWaP profiling to future UAV-suitable student models.The mentioned valuable related work will be appropriately cited. 
existing aerial datasets: The listed datasets are not directly usable for SSC evaluation, where occufly dataset outshines as they provide ground truth voxel grids. Few of the mentioned datasets especially the ones with depth maps can ofcource be benchmarked for depth estimation.

\paragraph{y7qY:}
The modest SSC mIoU suggests limitations beyond depth estimation: even the GT-depth oracle reaches only 7.25/8.40 mIoU, while OccuFly’s 96.109% empty voxels create severe occupancy imbalance, leaving only 3.891% to supervise occupied semantic classes.
The negative video result exposes limitations of our temporal adapter rather than temporal cues themselves, motivating motion/pose-aligned and occlusion-aware feature fusion used in recent temporal SSC.
The MVSFormer++-based FSSC gains are consistent with its cross-view fusion and attention-based cost-volume regularization providing a stronger geometric prior than our simple MVS module.
We have considered only scenes 1--8 for non standard  altitude interpolation ablation. 

\end{document}
